\documentclass[oneside]{book}
\usepackage[backend=biber,natbib=true,style=authoryear]{biblatex}
\addbibresource{/home/nguyen/1_NQBH/reference/bib.bib}
\usepackage[utf8]{inputenc}
\usepackage{float}
\usepackage{graphicx}
\usepackage[colorlinks=true,linkcolor=blue,urlcolor=red,citecolor=magenta]{hyperref}
\usepackage{amsmath,amssymb,amsthm}
\allowdisplaybreaks
\numberwithin{equation}{section}
\newtheorem{assumption}{Assumption}[section]
\newtheorem{lemma}{Lemma}[section]
\newtheorem{corollary}{Corollary}[section]
\newtheorem{definition}{Definition}[section]
\newtheorem{proposition}{Proposition}[section]
\newtheorem{theorem}{Theorem}[section]
\newtheorem{notation}{Notation}[section]
\newtheorem{remark}{Remark}[section]
\newtheorem{example}{Example}[section]
\newtheorem{ques}{Question}[section]
\newtheorem{problem}{Problem}[section]
\newtheorem{conjecture}{Conjecture}[section]
\usepackage[left=0.5in,right=0.5in,top=1.5cm,bottom=1.5cm]{geometry}
\usepackage{fancyhdr}
\pagestyle{fancy}
\fancyhf{}
\addtolength{\headheight}{0pt}% obsolete
\lhead{\scriptsize \chaptername~\thechapter}
\rhead{\scriptsize \nouppercase{\leftmark}} %\nouppercase !
\renewcommand{\chaptermark}[1]{\markboth{#1}{}}
\cfoot{\thepage}
\def\labelitemii{$\circ$}

\title{Psychology}
\author{Collector: Nguyen Quan Ba Hong}
\date{\today}

\begin{document}
\maketitle
\setcounter{secnumdepth}{6}
\tableofcontents

%------------------------------------------------------------------------------%

\chapter{Wikipedia's}

\section{\href{https://en.wikipedia.org/wiki/Psychology}{Wikipedia\texttt{/}Psychology}}


%------------------------------------------------------------------------------%

\section{\href{https://en.wikipedia.org/wiki/Narcissism}{Wikipedia\texttt{/}Narcissism}}
\textsf{\href{https://en.wikipedia.org/wiki/Narcissus_(Caravaggio)}{\textit{Narcissus}} (1597--99) by \href{https://en.wikipedia.org/wiki/Caravaggio}{Caravaggio}; the man in love with his own reflection.}

\textit{Narcissism} is the pursuit of gratification from \href{https://en.wikipedia.org/wiki/Vanity}{vanity} or \href{https://en.wikipedia.org/wiki/Egotistic}{egotistic} admiration of one's idealized \href{https://en.wikipedia.org/wiki/Self-image}{self-image} and attributes.

The term originated from \href{https://en.wikipedia.org/wiki/Greek_mythology}{Greek mythology}, where a young man named \href{https://en.wikipedia.org/wiki/Narcissus_(mythology)}{Narcissus} fell in love with his own image reflected in a pool of water.

Narcissism or pathological self-absorption was 1st identified as a disorder in 1898 by \href{https://en.wikipedia.org/wiki/Havelock_Ellis}{Havelock Ellis}[Millon, Theodore; Grossman, Seth; Million, Carrie; Meagher, Sarah; Ramnath, Rowena (2004). \textit{Personality Disorders in Modern Life}. Wiley. p. 343. ISBN 978-0-471-23734-1.] and featured in subsequent psychological models, e.g. in \href{https://en.wikipedia.org/wiki/Sigmund_Freud}{Freud}'s \href{https://en.wikipedia.org/wiki/On_Narcissism}{\textit{On Narcissism}} (1914).

The \href{https://en.wikipedia.org/wiki/American_Psychiatric_Association}{American Psychiatric Association} has listed the classification \href{https://en.wikipedia.org/wiki/Narcissistic_personality_disorder}{narcissistic personality disorder} in its \href{https://en.wikipedia.org/wiki/Diagnostic_and_Statistical_Manual_of_Mental_Disorders}{\textit{Diagnostic and Statistical Manual of Mental Disorders}} (DSM) since 1968, drawing on the historical concept of \href{https://en.wikipedia.org/wiki/Megalomania}{megalomania}.

%
It is distinct from concepts of distinguishing the self (\href{https://en.wikipedia.org/wiki/Egocentrism}{egocentrism} or \href{https://en.wikipedia.org/wiki/Egoism}{egoism}) and \href{https://en.wikipedia.org/wiki/Healthy_narcissism}{healthy forms of responsibility and care for oneself} (``\href{https://en.wikipedia.org/wiki/Primary_narcissism}{primary narcissism}'').

Narcissism by contrast is considered a problem for relationships with self and others and for maintaining a functional culture.

In \href{https://en.wikipedia.org/wiki/Trait_theory}{trait personality theory}, it features in several \href{https://en.wikipedia.org/wiki/Self-report_inventory}{self-report personality inventories} including the \href{https://en.wikipedia.org/wiki/Millon_Clinical_Multiaxial_Inventory}{Millon Clinical Multiaxial Inventory}.

It is 1 of the 3 \href{https://en.wikipedia.org/wiki/Dark_triad}{dark triadic} personality traits (the others being \href{https://en.wikipedia.org/wiki/Psychopathy}{psychopathy} and \href{https://en.wikipedia.org/wiki/Machiavellianism_(psychology)}{Machiavellianism}).

\subsection{History}

\subsection{Traits \& signs}

\subsection{Clinical \& research aspects}

\subsubsection{Narcissistic personality disorder}

\paragraph{Treatment \& management.}

\paragraph{Required element within normal development}

\paragraph{In relation to the pathological condition.}

\subsubsection{Commonly used measures}

\paragraph{Narcissistic Personality Inventory.}

\paragraph{Millon Clinical Multiaxial Inventory.}

\subsubsection{Empirical studies}

\subsubsection{Heritability research using twin studies}

\subsubsection{Stigmatizing attitude towards psychiatric illness}

\subsubsection{In evolutionary psychology}

\subsection{Narcissistic supply}

\subsection{Narcissistic rage \& narcissistic injury}

\subsection{Narcissistic defenses}

\subsection{Types}

\subsubsection{Masterson's subtypes (exhibitionist \& closet)}

\subsubsection{Millon's variations}

\subsubsection{Other forms}

\paragraph{Acquired situational narcissism.}

\paragraph{Codependency.}

\paragraph{Collective or group narcissism.}

\paragraph{Conversational narcissism.}

\paragraph{Cultural narcissism.}

\paragraph{Destructive narcissism.}

\paragraph{Malignant narcissism.}

\paragraph{Medical narcissism.}

\paragraph{In the workplace.}

\paragraph{Primordial narcissism.}

\paragraph{Sexual narcissism.}

\subsubsection{Narcissistic parents}

\subsubsection{Narcissistic leadership}

\subsection{In culture \& society}

\subsection{In fiction}

\subsection{See also}


%------------------------------------------------------------------------------%

\printbibliography[heading=bibintoc]
\end{document}